\section{Open problems}\label{sec:open}

\paragraph{The full bias boundary.}
Is $b_a$ strictly decreasing on $(a_*,1)$? It decreases near $a_*$
(Theorem~\ref{thm:bias-continuity}), and Figure~\ref{fig:phase}
suggests the same elsewhere. By \eqref{eq:bias-envelope-derivative},
this reduces to the sign at the unique optimizer. We do not know
that sign for all memories. The moment estimate gives
$b_a\le8(1-a)+O((1-a)^2)$. What is the exact decay near one?

\paragraph{A finite-length threshold.}
Does the central curvature have exactly one zero in $(0,1)$ at
every even length? Numerical zeros near $0.6563$ at $N=4000$ and
$0.6592$ at $N=8000$ suggest a slow increase toward $a_*$. Our limiting bias-interval theorem does not settle this
finite-length question.

\paragraph{The Drezner--Farnum model and higher dimensions.}
Changing the Drezner--Farnum innovation probability changes the
transition rule, so the reflected-mixture argument does not settle
its shapes when $p\ne1/2$~\cite{dreznerfarnum1993}.
In higher dimensions the fixed-point theory gives densities~\cite{glr2026},
but our sign argument requires points ordered along a line.
